\documentclass[11pt]{article}
\usepackage{amsmath,amssymb,amsthm}
\usepackage{fullpage}
\usepackage[hidelinks]{hyperref}
\newtheorem{theorem}{Theorem}
\newtheorem{lemma}{Lemma}
\theoremstyle{remark}
\newtheorem{remark}{Remark}
\newcommand{\todo}[1]{\textbf{[TODO: #1]}}

\title{The Dynamic Complementarity Tax:\\Two Consumers, One Record,
and an Aging Reference\\[4pt]\large Paper VII skeleton --- v0.1, for
venue work}
\author{A.~Bond \and S.~\todo{Syed's full name}}
\date{August 2026}

\begin{document}
\maketitle

\begin{abstract}
When two consumers read one stored record, the joint description
costs more than the better single-consumer description --- a
complementarity tax, on both the rate and the ideal reset-work faces.
We characterize how that tax responds when the eraser's reference
ages: every observation access class enters through one conditional
variance (equal-$q$ classes pay identical taxes --- a universality
that we prove is \emph{exactly} a Gaussian privilege and only
approximately a binary one); the tax's staleness response obeys an
envelope sign law (it rises iff the joint record couples to the
context harder than the binding consumer, with an always-downward
kink at the consumer switch); the whole theory lifts to jointly
stationary processes as per-frequency matrix water-level systems
under exactly two common Lagrange prices; and a new $m$-record
moment-convexity lemma makes the underlying weighted program convex,
resolving uniqueness and making the two-price allocation globally
optimal. A decode-threshold experiment measures the rising tax at
$+0.090$ against a $+0.086$ prediction, with an equal-$q$
analytic-equality control at $0.001$. \todo{compress per venue.}
\end{abstract}

\section{Introduction and related work}
\todo{Attribution spine from \texttt{go13-dynamic-tax-NOVELTY.md}:
Xiao--Luo 2005 (static kernel), Stylianou--Charalambous 2021/2024
(tuple reduction, implicit static two-water-level system ---
CRITICAL: also scopes the static substrate's rate face), HB/Kaspi +
Timo--Chan--Grant, vending-machine line. RE-SWEEP the
Stylianou--Charalambous line immediately pre-submission.}

\section{The static substrate}
\todo{Import GO-11 Thms 5/9 (statements only, cite the GO-11
manuscript); state the 074 uniqueness resolution here or in \S7.}

\section{The matrix-$q$ reduction and access-class universality}
\todo{Theorem 1 + the $r\ge2$ counterexample scoping.}

\section{The tax-curve sign law}
\todo{Theorem 2 (envelope; downward kink; $w{=}1$ flat; CT $\ge0$);
the exploratory regime map as a disclosed figure; the 067
reference-disambiguation story as a caution box.}

\section{The binary twin}
\todo{Theorem 3: slice collapse; reliability-distribution closed
form; generalized tilt; EXACT non-collapse at 50-digit precision vs
near-universality (envelope $\le2.1\times10^{-4}$, observation-grade)
--- the paper's most quotable dichotomy.}

\section{The spectral $m=2$ theorem}
\todo{Theorem 4: per-frequency Theorem-9 systems under two common
prices; the ``pair of prices'' vindication narrative (refuted at
$m=1$, exact at $m=2$); circulant scope + Toeplitz debts shared with
Paper VI.}

\section{The moment-convexity lemma}
\todo{Theorem 5 in full (four-step proof; $G=-\log\det(I-Z)$;
Hessian decomposition; flats; uniqueness scoping with the
shut-off-record failure mode). CANDIDATE: split out as a standalone
short communication if the venue prefers a leaner Paper VII.}

\section{The operational face}
\todo{069: instrument lineage, design frozen pre-instrument,
measured $+0.090$ vs $+0.086$, equal-$q$ control $0.001$; the V6
instrumentation amendment disclosed (both governed runs' physics
5/5).}

\section*{Reproducibility}
Sealed preregistrations 067--069, 072--074; governed harnesses
CI-re-run on every push; fresh-context adversarial verification for
every theorem (two draft refutations caught pre-seal, on record).
Map: \texttt{PAPER-VI-VII-EVIDENCE-INDEX.md}.

\todo{Bibliography per the attribution spine.}

\end{document}
